\section*{Capítulo \ref{ch:b3:residues} --- Séries de Laurent e
o teorema dos resíduos}

\begin{solution}{exo:b3:residues:1}
$\frac{\sin z}z = 1 - \frac{z^2}6 + \cdots$: removível,
\omterm{def:b3:residues:singularities}{resíduo} $0$. $\frac{\eu^z - 1}{z^2} = \frac1z + \frac12 +
\frac z6 + \cdots$: polo simples, \omterm{def:b3:residues:singularities}{resíduo} $1$.
$\frac1{z(z-1)^2}$: polo simples em $0$ com \omterm{def:b3:residues:singularities}{resíduo}
$\frac1{(0-1)^2} = 1$; polo duplo em $1$ com \omterm{def:b3:residues:singularities}{resíduo}
$\frac{\dd}{\dd z}\bigl(\frac1z\bigr)\big|_{z=1} = -1$.
$\frac{\cos z}{z^3} = \frac1{z^3} - \frac1{2z} + \cdots$: polo
de ordem $3$, \omterm{def:b3:residues:singularities}{resíduo} $-\frac12$. $\eu^{1/z} = \sum_{n\geq0}
\frac{z^{-n}}{n!}$: essencial, \omterm{def:b3:residues:singularities}{resíduo} $1$. $\frac1{\sin z}$:
polos simples em $n\pi$; \omterm{def:b3:residues:singularities}{resíduo} $\frac1{\cos 0} = 1$ em $0$,
$\frac1{\cos\pi} = -1$ em $\pi$
(\cref{met:b3:residues:compute}, $g/h'$).
\end{solution}

\begin{solution}{exo:b3:residues:2}
Com $z = \eu^{\iu t}$, $\cos t = \frac{z + z^{-1}}2$, $\dd t
= \frac{\dd z}{\iu z}$:
\[
\int_0^{2\pi}\frac{\dd t}{a + \cos t}
= \oint_{\abs z = 1}\frac{2\,\dd z}{\iu\,(z^2 + 2az + 1)} .
\]
As raízes $z_\pm = -a \pm \sqrt{a^2 - 1}$ satisfazem $z_+z_- =
1$ com $\abs{z_+} < 1 < \abs{z_-}$; o \omterm{def:b3:residues:singularities}{resíduo} em $z_+$ é
$\frac{1}{z_+ - z_-} = \frac1{2\sqrt{a^2-1}}$, de modo que a integral
vale $\frac2\iu\cdot2\iu\pi\cdot\frac1{2\sqrt{a^2-1}} =
\frac{2\pi}{\sqrt{a^2-1}}$. Quando $a \to 1^+$, ela explode (o
integrando tem pico em $t = \pi$); quando $a \to \infty$, ela se comporta
como $\frac{2\pi}a$, batendo com $\int\frac{\dd t}a$.
\end{solution}

\begin{solution}{exo:b3:residues:3}
Primeira integral: polos superiores de $\frac{z^2}{1+z^6}$ em $p =
\eu^{\iu\pi/6}, \iu, \eu^{5\iu\pi/6}$; em cada um,
$\operatorname{Res} = \frac{p^2}{6p^5} = \frac{p^3}{6p^6} =
-\frac{p^3}6$, e $p^3$ assume os valores $\iu, -\iu, \iu$:
soma dos \omterm{def:b3:residues:singularities}{resíduos} $-\frac{\iu}{6}$. O arco é
$O(R^{-4})\cdot O(R) \to 0$:
\[
\int_\R\frac{x^2\,\dd x}{1 + x^6} =
2\iu\pi\Bigl(-\frac \iu6\Bigr) = \frac\pi3 .
\]
Segunda: polo duplo em $\iu$ de $\frac1{(1+z^2)^2} =
\frac1{(z-\iu)^2(z+\iu)^2}$:
\[
\operatorname{Res} = \frac{\dd}{\dd z}\,(z +
\iu)^{-2}\Big|_{z=\iu} = \frac{-2}{(2\iu)^3} =
\frac{-2}{-8\iu} = -\frac\iu4,
\qquad
\int_\R\frac{\dd x}{(1+x^2)^2} =
2\iu\pi\cdot\Bigl(-\frac\iu4\Bigr) = \frac\pi2 ,
\]
coerente com o~\cref{exo:b3:fouriertransform:3}(b).
\end{solution}

\begin{solution}{exo:b3:residues:4}
Feche $\frac{\eu^{\iu tz}}{z^2 + a^2}$ no semiplano superior
($t \geq 0$): ali $\abs{\eu^{\iu tz}} =
\eu^{-t\operatorname{Im}z} \leq 1$, de modo que o arco contribui com
$O(R^{-2})\cdot O(R) \to 0$. O único polo encerrado $\iu a$
é simples, com \omterm{def:b3:residues:singularities}{resíduo} $\frac{\eu^{-at}}{2\iu a}$:
\[
\int_\R\frac{\eu^{\iu tx}}{x^2 + a^2}\dd x =
\frac{\pi}{a}\,\eu^{-at},
\qquad\text{logo}\qquad
\int_\R\frac{\cos(tx)}{x^2+a^2}\dd x = \frac\pi
a\,\eu^{-a\abs t}
\]
(parte real; par em $t$). Essa é a contrapartida por inversão
de $\widehat{\eu^{-a\abs x}} = \frac{2a}{a^2+\xi^2}$
(\cref{exo:b3:fouriertransform:1}): os dois cálculos
se confirmam mutuamente através do~\cref{thm:b3:fouriertransform:inversion}.
\end{solution}

\begin{solution}{exo:b3:residues:5}
Frações parciais: $f = \frac1{z-2} - \frac1{z-1}$.
Em $\abs z < 1$ (Taylor):
$f = \sum_{n\geq0}\bigl(1 - 2^{-n-1}\bigr)z^n$.
Em $1 < \abs z < 2$: $\frac1{z-2} =
-\sum_{n\geq0}\frac{z^n}{2^{n+1}}$ e $-\frac1{z-1} =
-\sum_{n\geq1}z^{-n}$: uma série bilateral genuína.
Em $\abs z > 2$: $f = \sum_{n\geq1}\bigl(2^{n-1} -
1\bigr)z^{-n}$. As três diferem porque as expansões de Laurent
estão ligadas a \emph{coroas}, não a pontos: cada região tem suas
próprias expansões geométricas. Os coeficientes de $z^{-1}$ ($-1$ e
$0$, respectivamente) são integrais sobre círculos que envolvem
\emph{as singularidades \omterm{def:b3:topology:interior}{interiores}}, não \omterm{def:b3:residues:singularities}{resíduos} em $0$ ($f$ é
\omterm{def:b3:holomorphic:holo}{holomorfa} em $0$): para $1 < \abs z < 2$ o coeficiente
$-1$ é $\operatorname{Res}(f, 1)$; para $\abs z > 2$ o
coeficiente $0$ é $\operatorname{Res}(f,1) +
\operatorname{Res}(f,2) = -1 + 1$. Os \omterm{def:b3:residues:singularities}{resíduos} de $f$: $-1$
em $1$ e $+1$ em $2$.
\end{solution}

\begin{solution}{exo:b3:residues:6}
(a) A \omterm{thm:b3:residues:laurent}{série de Laurent} $\sum_nz^{-n}/n!$ tem infinitos termos
negativos: essencial. Resolvendo $\eu^{1/z} = w$ ($w \neq
0$): $\frac1z = \log\abs w + \iu\arg w + 2\iu\pi k$, de modo que
\[
z_k = \frac1{\log\abs w + \iu\arg w + 2\iu\pi k}
\xrightarrow[k\to\infty]{} 0 :
\]
todo valor não nulo é atingido uma infinidade de vezes perto de $0$ ---
mais forte que \omterm{ex:b3:lebesgue:gamma}{densidade}.
(b) Ao longo de $z = 1/x$, $x \to +\infty$: $\eu^x \to \infty$;
ao longo de $z = \iu/y$: módulo $1$. Um polo exige
$\abs{f(z)} \to \infty$ \emph{ao longo de toda aproximação}: aqui o
limite simplesmente não existe, nem mesmo em $[0, +\infty]$.
\end{solution}

\begin{solution}{exo:b3:residues:7}
(a) Em $\abs z = 1$: $\abs{z^7 + z - 1} \leq 3 < 4 =
\abs{-4z^3}$. Rouché com $f = -4z^3$, $g = z^7 + z - 1$:
três zeros no disco.
(b) Em $\abs z = 1$: $\abs{z^4 + 1} \leq 2 < 5 = \abs{5z}$:
um zero em $\abs z < 1$. Em $\abs z = 2$: $\abs{5z + 1} \leq
11 < 16 = \abs{z^4}$: quatro zeros em $\abs z < 2$. Logo, três
zeros na coroa.
(c) Para $P = z^n + a_{n-1}z^{n-1} + \dots$: em $\abs z = R >
1 + \sum\abs{a_k}$, $\abs{P - z^n} \leq
\bigl(\sum\abs{a_k}\bigr)R^{n-1} < R^n = \abs{z^n}$: $P$ tem
exatamente $n$ zeros em $D(0, R)$ --- d'Alembert--Gauss com
multiplicidade, por pura contagem.
\end{solution}

\begin{solution}{exo:b3:residues:8}
(a) Suponha $f(a) = 0$, $f \not\equiv 0$: escolha $r$ com $f$
sem zeros no círculo $C = \partial D(a, r)$ (\omterm{thm:b3:holomorphic:identity}{zeros
isolados}) e $m = \min_C\abs f > 0$. Pelo~\cref{thm:b3:holomorphic:weierstrassconv}, $f_n \to f$ e
$f_n' \to f'$ uniformemente em $C$; para $n$ grande, $\abs{f_n}
\geq m/2$ em $C$, de modo que
\[
\frac1{2\iu\pi}\int_C\frac{f_n'}{f_n}
\longrightarrow \frac1{2\iu\pi}\int_C\frac{f'}{f} \geq 1
\]
(o limite conta o zero $a$; a convergência porque
os numeradores convergem uniformemente e os denominadores são uniformemente
minorados). O membro esquerdo é um inteiro que conta zeros de
$f_n$ no disco: ele tem de ser $\geq 1$ a partir de certa ordem ---
contradizendo a ausência de zeros. Logo $f$ não tem zeros.

(b) Fixe $w \in \Omega$ e aplique (a) no aberto \omterm{def:b3:topology:connected}{conexo}
$\Omega\setminus\{w\}$ (remover um ponto de um subconjunto aberto
\omterm{def:b3:topology:connected}{conexo} de $\C$ preserva a \omterm{def:b3:topology:connected}{conexidade}) a $g_n(z) = f_n(z) -
f_n(w)$, sem zeros ali por injetividade, convergindo para $g = f
- f(w)$. Se $f$ não é constante, $g \not\equiv 0$ em
$\Omega\setminus\{w\}$, de modo que $g$ não tem zeros ali: $f(z) \neq
f(w)$ para todo $z \neq w$. Como $w$ era arbitrário, $f$ é
injetora.
\end{solution}

\begin{solution}{exo:b3:residues:9}
O bordo do setor consiste em $[0, R]$, no arco $A_R$ e na
semirreta $\eu^{2\iu\pi/n}[0, R]$ percorrida ao contrário. Dentro está o
único polo $p = \eu^{\iu\pi/n}$ de $\frac1{1+z^n}$, com
\omterm{def:b3:residues:singularities}{resíduo} $\frac1{np^{n-1}} = \frac{p}{np^n} = -\frac pn$. Na
semirreta de volta, $z = \eu^{2\iu\pi/n}x$ dá $z^n = x^n$ e
$\dd z = \eu^{2\iu\pi/n}\dd x$; o arco é
$O(R^{-n})\cdot O(R) \to 0$. Logo
\[
\bigl(1 - \eu^{2\iu\pi/n}\bigr)
\int_0^\infty\frac{\dd x}{1 + x^n}
= 2\iu\pi\Bigl(-\frac{\eu^{\iu\pi/n}}n\Bigr),
\quad\text{logo}\quad
\int_0^\infty\frac{\dd x}{1+x^n}
= \frac{2\iu\pi}{n\,\bigl(\eu^{\iu\pi/n} -
\eu^{-\iu\pi/n}\bigr)} = \frac{\pi}{n\sin(\pi/n)} .
\]
$n = 2$: $\frac\pi{2\sin(\pi/2)} = \frac\pi2 =
[\arctan]_0^\infty$. Quando $n \to \infty$: o valor tende a
$1$ e, de fato, o integrando tende a $\mathbf 1_{\intco01}$
(com dominação para o TCD $\min(1, x^{-2})$ para $n \geq 2$).
\end{solution}

\begin{solution}{exo:b3:residues:10}
Em $\abs z = R$ grande, $\abs f \leq C R^{-2}$: $\abs{\oint_{C_R}
f} \leq 2\pi R\cdot CR^{-2} \to 0$. Mas, para $R$ além de todos os
polos, o teorema dos \omterm{def:b3:residues:singularities}{resíduos} dá $\oint_{C_R}f =
2\iu\pi\sum_{\text{todo }p}\operatorname{Res}(f, p)$: a soma
total se anula. Para $f = \frac1{z(z-1)(z-2)}$: \omterm{def:b3:residues:singularities}{resíduos}
$\frac1{(-1)(-2)} = \frac12$ em $0$, $\frac1{1\cdot(-1)} = -1$
em $1$, $\frac1{2\cdot1} = \frac12$ em $2$ --- somando $0$,
como previsto, e
\[
\frac1{z(z-1)(z-2)} = \frac{1/2}{z} - \frac1{z - 1} +
\frac{1/2}{z-2} :
\]
os \omterm{def:b3:residues:singularities}{resíduos} \emph{são} os coeficientes das frações parciais, e
a identidade de soma nula fornece uma verificação de coerência gratuita (ou
determina o último coeficiente a partir dos demais).
\end{solution}

\begin{solution}{exo:b3:residues:11}
No buraco de fechadura, com a determinação escolhida:
logo acima do corte, $\log z = \ln x$; logo abaixo, $\log z =
\ln x + 2\iu\pi$. As quatro peças dão
\[
\Bigl(1 - \eu^{2\iu\pi(a-1)}\Bigr)\int_\varepsilon^R
\frac{x^{a-1}}{1+x}\dd x + \int_{C_R} + \int_{C_\varepsilon}
= 2\iu\pi\operatorname{Res}(f, -1) .
\]
Arcos: $\abs{f} \leq \frac{R^{a-1}}{R - 1}$ em $C_R$, de comprimento
$2\pi R$: contribuição $O(R^{a-1}) \to 0$ ($a < 1$);
$\abs f \leq \frac{\varepsilon^{a-1}}{1 - \varepsilon}$ em
$C_\varepsilon$, de comprimento $2\pi\varepsilon$: $O(\varepsilon^a)
\to 0$ ($a > 0$). \omterm{def:b3:residues:singularities}{Resíduo}: em $z = -1 = \eu^{\iu\pi}$,
$\operatorname{Res} = \eu^{(a-1)\iu\pi}$. Logo
\[
I(a) = \frac{2\iu\pi\,\eu^{\iu\pi(a-1)}}{1 -
\eu^{2\iu\pi(a-1)}}
= \frac{2\iu\pi}{\eu^{-\iu\pi(a-1)} - \eu^{\iu\pi(a-1)}}
= \frac{\pi}{-\sin(\pi(a-1))} = \frac{\pi}{\sin\pi a} .
\]
Reflexão da gama: $B(a, 1-a) =
\int_0^1t^{a-1}(1-t)^{-a}\dd t$; a substituição $t =
\frac x{1+x}$, $1 - t = \frac1{1+x}$, $\dd t =
\frac{\dd x}{(1+x)^2}$ a transforma em $\int_0^\infty
\frac{x^{a-1}}{1+x}\dd x = I(a)$, e a fórmula de Euler
$B(a, 1-a) = \Gamma(a)\Gamma(1-a)/\Gamma(1) $
(\cref{pb:b3:lebesgue:1}) dá
$\Gamma(a)\Gamma(1-a) = \frac\pi{\sin\pi a}$ --- em
particular, $\Gamma(\tfrac12) = \sqrt\pi$ mais uma vez.

\end{solution}

\begin{solution}{exo:b3:residues:12}
(a) Em $\abs z = 1$: $\abs{z^4} = 1 < 7 \leq \abs{8z + 1}$
($\abs{8z} - 1 = 7$): $P$ tem tantos zeros em $\mathbb D$
quantos $8z + 1$, a saber, \emph{um} (em $-\frac18$).

(b) Em $\abs z = 3$: $\abs{8z + 1} \leq 25 < 81 =
\abs{z^4}$: Rouché contra $z^4$ dá os quatro zeros em
$\abs z < 3$, logo $4 - 1 = 3$ zeros na coroa $1 <
\abs z < 3$. Refinamento: um zero com $\abs z = r \geq 2.1$
satisfaria $r^4 = \abs{8z + 1} \leq 8r + 1$, mas $r^4 -
8r - 1$ é crescente para $r \geq 2$ e vale $19.45 -
16.8 - 1 = 1.65 > 0$ em $r = 2.1$: impossível. Logo os três
zeros externos estão em $1 < \abs z < 2.1$. (Numericamente: um zero real
perto de $-1.95$ e um par conjugado perto de $1.04 \pm
1.73\iu$, de módulo $2.02$ --- razão pela qual uma tentativa de Rouché
no raio exatamente $2$ tem de falhar: o teorema exige
dominação estrita, e os zeros ficam logo do lado de fora.)

(c) Sem zeros em $\iu\R$: $\operatorname{Re}P(\iu t) = t^4 +
1 \geq 1$. Zeros no semiplano direito: use o princípio do
argumento no bordo do semidisco $\{\abs z \leq
R,\ \operatorname{Re}z \geq 0\}$. No arco grande, $\arg P
\approx \arg z^4$ gira de $4\cdot\pi = 2\pi\cdot2$ (o
arco abarca o ângulo $\pi$). Ao longo do eixo imaginário, de
$\iu R$ até $-\iu R$: $P(\iu t) = (t^4 + 1) + 8\iu t$
permanece no semiplano direito ($\operatorname{Re} > 0$), de modo que
$\arg P$ varia dentro de $\intoo{-\pi/2}{\pi/2}$ e retorna
com variação líquida $\to 0$ quando $R \to \infty$ (extremos ambos
$\approx \arg t^4 = 0$). Enrolamento total: $\frac{4\pi + 0}
{2\pi} = 2$: \emph{dois} zeros no semiplano direito ---
coerente com a numérica: o par conjugado $\approx
1.04 \pm 1.73\iu$ tem parte real positiva, e os zeros reais
$\approx -0.125$ e $\approx -1.96$, negativa.
\end{solution}

\begin{solution}{pb:b3:residues:1}
\textbf{1.} $\sin(\pi z)$ tem zeros simples exatamente em $\Z$
($\sin\pi z = 0$ se, e somente se, $z \in \Z$, e $(\sin\pi z)' = \pi\cos\pi
z \neq 0$ ali), e $\cos(\pi n) \neq 0$: $\pi\cot(\pi z) =
\pi\cos(\pi z)/\sin(\pi z)$ tem polos simples em $\Z$ com
\[
\operatorname{Res}(\pi\cot\pi z,\ n)
= \frac{\pi\cos(\pi n)}{\pi\cos(\pi n)} = 1
\]
(regra $g/h'$, \cref{met:b3:residues:compute}).

\textbf{2.} $\frac{\sin(\pi z)}{\pi z} = 1 - \frac{(\pi z)^2}6
+ \frac{(\pi z)^4}{120} - \cdots$ é \omterm{def:b3:holomorphic:holo}{holomorfa} e não nula
perto de $0$: seu recíproco é \omterm{def:b3:holomorphic:holo}{holomorfo}
(\cref{def:b3:holomorphic:holo}: quociente), com série $1 +
\frac{(\pi z)^2}{6} + \frac{7(\pi z)^4}{360} + \cdots$
(identifique: coeficientes ao estilo $(1 - u)^{-1}$ a partir de $u =
\frac{(\pi z)^2}6 - \frac{(\pi z)^4}{120}$: o coeficiente de $z^4$
é $\frac1{36} - \frac1{120} = \frac{7}{360}$).
Multiplique por $\cos(\pi z) = 1 - \frac{(\pi z)^2}2 + \frac{(\pi
z)^4}{24} - \cdots$ e divida por $z$:
\[
\pi\cot(\pi z) = \frac1z\Bigl[1 + \pi^2z^2\Bigl(\frac16 -
\frac12\Bigr) + \pi^4z^4\Bigl(\frac7{360} - \frac1{12} +
\frac1{24}\Bigr)\Bigr] + \cdots
= \frac1z - \frac{\pi^2}3\,z - \frac{\pi^4}{45}\,z^3 -
\cdots
\]
($\frac7{360} - \frac{30}{360} + \frac{15}{360} =
-\frac8{360} = -\frac1{45}$).

\textbf{3.} Lados verticais $z = \pm(N + \frac12) + \iu y$:
pela $\pi$-periodicidade de $\cot$, $\cot(\pi z) = \cot(\pm\frac\pi2
+ \iu\pi y) = -\tan(\iu\pi y) = -\iu\tanh(\pi y)$, de módulo
$\leq 1$. Lados horizontais $z = x \pm \iu(N + \frac12)$: de
$\abs{\cot(a + \iu b)}^2 = \frac{\cos^2a +
\sinh^2b}{\sin^2a + \sinh^2b} \leq \frac{1 +
\sinh^2b}{\sinh^2b} = \coth^2 b$,
\[
\abs{\cot(\pi z)} \leq \coth\bigl(\pi(N + \tfrac12)\bigr)
\leq \coth(\pi/2) < 1.1 .
\]
Ambas as cotas são $\leq 2$.

\textbf{4.} Aplique o~\cref{thm:b3:residues:residue} a $F(z) =
\pi\cot(\pi z)f(z)$ em $C_N$ ($N$ além de todos os polos de $f$):
\[
\frac1{2\iu\pi}\oint_{C_N}F = \sum_{n=-N}^{N}f(n) +
\sum_p\operatorname{Res}(F, p),
\]
contribuindo os polos inteiros com $f(n)$ (questão 1; $f$
\omterm{def:b3:holomorphic:holo}{holomorfa} ali). Em $C_N$: $\abs F \leq 2\pi\cdot
C\abs z^{-2} \leq 2\pi C N^{-2}$, e o perímetro é $8(N +
\frac12)$: a integral é $O(1/N) \to 0$. Faça $N \to
\infty$: a fórmula de soma em destaque.

\textbf{5.} O decaimento $\abs f = O(\abs z^{-2})$ matou a
\omterm{def:b3:holomorphic:contour}{integral de contorno} \emph{e} fez $\sum\abs{f(n)}$ convergir.
Para $f(z) = \frac1{z + \frac12}$: as somas simétricas
$\sum_{-N}^N\frac1{n + \frac12}$ telescopam para $0$ (os termos
$n$ e $-n - 1$ se cancelam), e o membro direito é
$-\operatorname{Res}\bigl(\frac{\pi\cot\pi z}{z + \frac12},
-\frac12\bigr) = -\pi\cot(-\frac\pi2) = 0$: coerente ---
mas $\sum\abs{f(n)}$ diverge; o método calcula apenas o
limite simétrico (valor principal).

\textbf{6.} $g(z) = \frac{\pi\cot(\pi z)}{z^2}$: polos nos
inteiros não nulos com \omterm{def:b3:residues:singularities}{resíduos} $\frac1{n^2}$, e em $0$,
onde, pela questão 2,
\[
g(z) = \frac1{z^3} - \frac{\pi^2}{3z} - \frac{\pi^4}{45}z -
\cdots :
\qquad \operatorname{Res}(g, 0) = -\frac{\pi^2}3 .
\]
A \omterm{def:b3:holomorphic:contour}{integral de contorno} sobre $C_N$ tende a $0$ como na questão 4
($\abs{g} = O(N^{-2})$ em $C_N$). Logo $0 =
\sum_{n\neq0}\frac1{n^2} - \frac{\pi^2}3$: $2\zeta(2) =
\frac{\pi^2}3$, $\zeta(2) = \frac{\pi^2}6$.

\textbf{7.} $g(z) = \frac{\pi\cot(\pi z)}{z^4} = \frac1{z^5}
- \frac{\pi^2}{3z^3} - \frac{\pi^4}{45z} - \cdots$: \omterm{def:b3:residues:singularities}{resíduo}
em $0$ igual a $-\frac{\pi^4}{45}$, e $0 =
2\zeta(4) - \frac{\pi^4}{45}$: $\zeta(4) =
\frac{\pi^4}{90}$.

\textbf{8.} Com $g = \pi\cot(\pi z)/z^{2k}$: o \omterm{def:b3:residues:singularities}{resíduo} em
$0$ é o coeficiente $a_{2k-1}$ de $z^{2k-1}$ na
expansão de $\pi\cot(\pi z)$, e a anulação do contorno dá
$2\zeta(2k) + a_{2k-1} = 0$: $\zeta(2k) = -a_{2k-1}/2$, um
múltiplo racional de $\pi^{2k}$, já que os coeficientes da cotangente
o são. Mais um passo de divisão fornece $a_5 =
-\frac{2\pi^6}{945}$, donde $\zeta(6) = \frac{\pi^6}{945}$.
Para $\zeta(3)$: o núcleo natural $g = \pi\cot(\pi z)/z^3$
produz $\sum_{n\neq0}\frac1{n^3} = 0$ por imparidade --- o
método demonstra $0 = 0$ e é estruturalmente cego aos valores
ímpares da zeta (nenhuma forma fechada para $\zeta(3)$ é conhecida; sua
irracionalidade, Apéry 1978, exigiu ideias inteiramente
diferentes).

\textbf{9.} $F(z) = \frac{\pi\cot(\pi z)}{(z - w)(z + w)}$
tem polos nos inteiros (\omterm{def:b3:residues:singularities}{resíduos} $\frac1{n^2 - w^2}$,
notando o sinal: $f(n) = \frac1{(n-w)(n+w)} = \frac1{n^2 -
w^2}$) e polos simples em $\pm w$ com \omterm{def:b3:residues:singularities}{resíduos}
$\frac{\pi\cot(\pm\pi w)}{\pm2w} = \frac{\pi\cot(\pi
w)}{2w}$ cada ($\cot$ é ímpar). O argumento da questão 4
($\abs f = O(\abs z^{-2})$) dá
\[
\sum_{n\in\Z}\frac1{n^2 - w^2} + \frac{\pi\cot(\pi w)}{w} =
0,
\qquad\text{isto é}\qquad
\pi\cot(\pi w) = \frac1w + \sum_{n\geq1}\frac{2w}{w^2 -
n^2},
\]
(o termo $n = 0$ é $-\frac1{w^2}$; reagrupe $\pm n$). Convergência
normal nos \omterm{def:b3:topology:compact}{compactos} de $\C\setminus\Z$: para $\abs w \leq
R$ e $n \geq 2R$, $\abs{\frac{2w}{w^2 - n^2}} \leq
\frac{2R}{n^2 - R^2} \leq \frac{8R}{3n^2}$.

\textbf{10.} Para $\abs w \leq r < 1$: $\frac{2w}{w^2 - n^2} =
-\frac{2w}{n^2}\cdot\frac1{1 - w^2/n^2} =
-2\sum_{k\geq0}\frac{w^{2k+1}}{n^{2k+2}}$, com
$\abs{\text{termos}} \leq 2r^{2k+1}/n^{2k+2}$, somável em
$(n, k)$: Fubini para séries rearranja
\[
\pi\cot(\pi w) = \frac1w -
2\sum_{k\geq0}\zeta(2k+2)\,w^{2k+1} .
\]
Comparando com a questão 2: $-2\zeta(2) = -\frac{\pi^2}3$ e
$-2\zeta(4) = -\frac{\pi^4}{45}$ --- os mesmos valores. Três
\omterm{def:b3:topology:pathconnected}{caminhos} para $\frac{\pi^2}6$: \omterm{thm:b3:hilbert:parseval}{Parseval} (séries de Fourier), o
traço do operador de Green da corda e as duas expansões da
cotangente; que uma soma sobre frequências, um traço de operador
e uma \omterm{def:b3:holomorphic:contour}{integral de contorno} concordem não é acidente --- cada um é uma
face da mesma identidade espectral.

\textbf{11.} Fixe $R \geq 1$ e seja $n_0$ o menor
inteiro $\geq 2R$. Para $\abs z \leq R$ e $n \geq n_0$:
$\abs{z^2/n^2} \leq \frac14$, de modo que $1 - z^2/n^2 \in
\bar D(1, \frac14)$, em que o logaritmo principal é
\omterm{def:b3:holomorphic:holo}{holomorfo}, e
\[
\Bigl|\log\Bigl(1 - \frac{z^2}{n^2}\Bigr)\Bigr|
\leq \sum_{k\geq1}\frac1k\,\Bigl|\frac{z^2}{n^2}\Bigr|^k
\leq \frac{\abs{z^2/n^2}}{1 - \abs{z^2/n^2}}
\leq \frac{2R^2}{n^2} :
\]
a soma $S(z) = \sum_{n\geq n_0}\log(1 - z^2/n^2)$ converge
normalmente em $\bar D(0, R)$, com somas parciais \omterm{def:b3:holomorphic:holo}{holomorfas}
$S_N$ e $\abs{S_N} \leq 2R^2\zeta(2)$ uniformemente. Como
$\abs{\eu^a - \eu^b} \leq
\eu^{\max(\abs a,\abs b)}\abs{a - b}$ (cota do valor médio no
segmento), os produtos das caudas $\prod_{n_0\leq n\leq N} =
\eu^{S_N}$ convergem uniformemente em $\bar D(0, R)$ para a
função sem zeros $\eu^S$. Multiplicando pelo polinômio fixo
$z\prod_{n<n_0}(1 - z^2/n^2)$: $P_N \to P$ uniformemente em
$\bar D(0, R)$, e o~\cref{thm:b3:holomorphic:weierstrassconv} torna $P$
\omterm{def:b3:holomorphic:holo}{holomorfa} ali; sendo $R$ arbitrário, $P$ é inteira. Em
$\bar D(0, R)$, os zeros de $P$ são os do prefator
polinomial --- os inteiros de módulo $\leq R$, cada um simples
($(1 - z/n)(1 + z/n)$ tem zeros simples distintos, $\eu^S$
nenhum): o conjunto de zeros de $P$ é $\Z$, todos simples.

\textbf{12.} Derivação logarítmica do produto finito,
fora de seus zeros:
\[
\frac{P_N'(z)}{P_N(z)} = \frac1z +
\sum_{n=1}^{N}\frac{-2z/n^2}{1 - z^2/n^2}
= \frac1z + \sum_{n=1}^{N}\frac{2z}{z^2 - n^2} .
\]
Num \omterm{def:b3:topology:compact}{compacto} $K \subseteq \C\setminus\Z$: $P_N \to P$ e
$P_N' \to P'$ uniformemente
(\cref{thm:b3:holomorphic:weierstrassconv}), e $\min_K
\abs P > 0$ ($P$ só se anula em $\Z$), de modo que, a partir de certa ordem,
$\abs{P_N} \geq \frac12\min_K\abs P$ e $P_N'/P_N \to
P'/P$ uniformemente em $K$. O membro do meio converge para
$\frac1z + \sum_{n\geq1}\frac{2z}{z^2-n^2} = \pi\cot(\pi z)$
pela questão 9: logo $P'/P = \pi\cot(\pi z)$ em
$\C\setminus\Z$.

\textbf{13.} $\sin(\pi z)$ e $P$ são inteiras com o mesmo
conjunto de zeros $\Z$, todos simples (questões 1 e 11). Perto de
$m \in \Z$, escreva $\sin(\pi z) = (z - m)\,\sigma(z)$ e
$P(z) = (z - m)\,\psi(z)$ com $\sigma, \psi$ \omterm{def:b3:holomorphic:holo}{holomorfas}
e não nulas em $m$ (fatore a série de potências): $Q =
\sin(\pi z)/P = \sigma/\psi$ se estende \omterm{def:b3:holomorphic:holo}{holomorficamente} e
sem zeros através de cada inteiro, e é sem zeros em
$\C\setminus\Z$ como quociente de funções sem zeros. Ali,
\[
\frac{Q'}{Q} = \frac{(\sin\pi z)'}{\sin\pi z} -
\frac{P'}{P} = \pi\cot(\pi z) - \pi\cot(\pi z) = 0 ,
\]
de modo que a função inteira $Q'$ se anula em $\C\setminus\Z$,
logo em toda parte por \omterm{def:b3:topology:continuity}{continuidade}: $Q$ é constante. Quando $z \to
0$: $\sin(\pi z)/z \to \pi$ e $P(z)/z \to 1$, de modo que $Q =
\pi$:
\[
\sin(\pi z) = \pi z\prod_{n\geq1}
\Bigl(1 - \frac{z^2}{n^2}\Bigr) .
\]

\textbf{14.} Em $z = \frac12$: $1 = \sin\frac\pi2 =
\frac\pi2\prod_{n\geq1}\bigl(1 - \frac1{4n^2}\bigr) =
\frac\pi2\prod\frac{4n^2-1}{4n^2}$, de modo que
\[
\frac\pi2 = \prod_{n\geq1}\frac{4n^2}{4n^2 - 1}
= \lim_{N\to\infty}\prod_{n=1}^N
\frac{(2n)(2n)}{(2n-1)(2n+1)}
= \lim_{N\to\infty}
\frac{2\cdot2\cdot4\cdot4\cdots(2N)(2N)}
{1\cdot3\cdot3\cdot5\cdots(2N-1)(2N+1)} :
\]
o produto de Wallis, um corolário de uma linha da fatoração de
Euler.

\textbf{15.} Para $\abs z \leq r < 1$, todo fator está em
$D(1, r^2) \subseteq D(1, 1)$, de modo que $P(z)/z = \exp\bigl(
\sum_{n\geq1}\log(1 - z^2/n^2)\bigr)$: cada produto parcial
é a exponencial de uma soma parcial, e ambos os membros passam
ao limite pela \omterm{def:b3:topology:continuity}{continuidade} de $\exp$. A série dupla
\[
\sum_{n\geq1}\log\Bigl(1 - \frac{z^2}{n^2}\Bigr)
= -\sum_{n\geq1}\sum_{k\geq1}\frac{z^{2k}}{k\,n^{2k}}
= -\sum_{k\geq1}\frac{\zeta(2k)}k\,z^{2k}
= -\zeta(2)\,z^2 + O(z^4)
\]
se rearranja por Fubini para séries: $\sum_{n,k}
\frac{r^{2k}}{kn^{2k}} \leq \sum_k\zeta(2k)r^{2k} \leq
\zeta(2)\frac{r^2}{1-r^2} < \infty$. Logo
\[
P(z) = z\,\exp\bigl(-\zeta(2)z^2 + O(z^4)\bigr)
= z - \zeta(2)\,z^3 + O(z^5) ,
\]
e a questão 13 compara isso com $\sin(\pi z) = \pi z -
\frac{\pi^3}6z^3 + O(z^5)$: $\pi\zeta(2) = \frac{\pi^3}6$,
isto é, $\zeta(2) = \frac{\pi^2}6$. A face aditiva
(questão 10) e a face multiplicativa calculam o mesmo
número.

\textbf{16.} Num \omterm{def:b3:topology:compact}{compacto} $K \subseteq \C\setminus\Z$, as
somas parciais $S_N = \frac1z + \sum_{n\leq N}\bigl(
\frac1{z-n} + \frac1{z+n}\bigr)$ (questão 9, com os termos
reagrupados como $\frac{2z}{z^2-n^2} = \frac1{z-n} +
\frac1{z+n}$) convergem uniformemente para $\pi\cot(\pi z)$, de modo que
o~\cref{thm:b3:holomorphic:weierstrassconv} dá $S_N' \to
(\pi\cot\pi z)' = -\pi^2/\sin^2(\pi z)$ uniformemente em $K$.
Como $S_N' = -\sum_{\abs n\leq N}(z - n)^{-2}$:
\[
\frac{\pi^2}{\sin^2(\pi z)} =
\sum_{n\in\Z}\frac1{(z - n)^2} ,
\]
sendo a convergência normal nos \omterm{def:b3:topology:compact}{compactos} de $\C\setminus\Z$
(termos $O(n^{-2})$).

\textbf{17.} Em $z = \frac12$, o membro esquerdo é $\pi^2$; à
direita, $(\frac12 - n)^2 = \frac{(2n-1)^2}4$, com $2n -
1$ percorrendo todos os inteiros ímpares exatamente uma vez quando $n$ percorre
$\Z$:
\[
\pi^2 = \sum_{n\in\Z}\frac{4}{(2n-1)^2}
= 8\sum_{m\geq0}\frac1{(2m+1)^2},
\qquad
\sum_{m\geq0}\frac1{(2m+1)^2} = \frac{\pi^2}8 .
\]
Separação por paridade: $\zeta(2) = \frac{\pi^2}8 +
\sum_{n\geq1}\frac1{(2n)^2} = \frac{\pi^2}8 +
\frac{\zeta(2)}4$, de modo que $\frac34\zeta(2) = \frac{\pi^2}8$ e
$\zeta(2) = \frac{\pi^2}6$ mais uma vez.

\textbf{18.} Com $c = \cot\theta$ e $\cot(2\theta) =
\frac{c^2-1}{2c}$: $\cot\theta - 2\cot(2\theta) = c -
\frac{c^2-1}c = \frac1c = \tan\theta$. Logo $\pi\tan(\pi z)
= \pi\cot(\pi z) - 2\pi\cot(2\pi z)$, e a questão 9 em $z$
e em $2z$ dá
\[
\pi\cot(\pi z) = \frac1z +
\sum_{n\geq1}\frac{2z}{z^2 - n^2},
\qquad
2\pi\cot(2\pi z) = \frac1z +
\sum_{n\geq1}\frac{8z}{4z^2 - n^2} .
\]
Ambas as séries convergem absolutamente em cada $z$ fixo fora dos
polos, de modo que a diferença pode ser reagrupada à vontade: na
segunda série, os termos pares $n = 2m$ dão $\frac{8z}{4z^2 -
4m^2} = \frac{2z}{z^2 - m^2}$ e cancelam inteiramente a primeira série,
restando
\[
\pi\tan(\pi z) = -\sum_{m\geq0}\frac{8z}{4z^2 - (2m+1)^2}
= \sum_{m\geq0}\frac{8z}{(2m+1)^2 - 4z^2} ,
\]
normalmente nos \omterm{def:b3:topology:compact}{compactos} que evitam $\frac12 + \Z$ (termos
$O(m^{-2})$).

\textbf{19.} Para $\abs z \leq r < \frac12$:
\[
\frac{8z}{(2m+1)^2 - 4z^2} = \frac{8z}{(2m+1)^2}
\sum_{k\geq0}\Bigl(\frac{4z^2}{(2m+1)^2}\Bigr)^{k},
\]
com $\sum_{m,k}8r\,(4r^2)^k(2m+1)^{-2k-2} < \infty$, pois
$4r^2 < 1$: Fubini rearranja a soma dupla em
\[
\pi\tan(\pi z) =
\sum_{k\geq0}8\cdot4^k\,\lambda(2k+2)\,z^{2k+1} .
\]
Contra $\pi\tan(\pi z) = \pi^2z + \frac{\pi^4}3z^3 +
O(z^5)$: o coeficiente de $z$ dá $8\lambda(2) = \pi^2$
--- de novo a questão 17 --- e o de $z^3$ dá
$32\lambda(4) = \frac{\pi^4}3$, isto é, $\lambda(4) =
\frac{\pi^4}{96}$. Removendo os denominadores pares,
$\lambda(4) = \zeta(4) - 2^{-4}\zeta(4) =
\frac{15}{16}\zeta(4)$: $\zeta(4) = \frac{16}{15}\cdot
\frac{\pi^4}{96} = \frac{\pi^4}{90}$, batendo com a questão 7.

\textbf{20.} $\cot\frac\theta2 - \cot\theta =
\frac{\cos\frac\theta2\,\sin\theta -
\cos\theta\,\sin\frac\theta2}{\sin\frac\theta2\,\sin\theta}
= \frac{\sin\frac\theta2}{\sin\frac\theta2\,\sin\theta} =
\frac1{\sin\theta}$, sendo o numerador $\sin(\theta -
\frac\theta2)$. Com $\theta = \pi z$, a questão 9 em
$\frac z2$ se lê $\pi\cot\frac{\pi z}2 = \frac2z +
\sum_{n\geq1}\frac{4z}{z^2 - 4n^2}$, de modo que
\[
\frac{\pi}{\sin(\pi z)}
= \pi\cot\frac{\pi z}2 - \pi\cot(\pi z)
= \frac1z + \sum_{n\geq1}\frac{4z}{z^2 - 4n^2}
- \sum_{n\geq1}\frac{2z}{z^2 - n^2} .
\]
A convergência absoluta permite o reagrupamento por paridade: os
$n
= 2m$ pares na série subtraída contribuem com
$\frac{2z}{z^2-4m^2}$, restando $+\frac{2z}{z^2-4m^2}$ da
primeira soma, ao passo que os $n$ ímpares sobrevivem com sinal $-$:
\[
\frac{\pi}{\sin(\pi z)} = \frac1z +
\sum_{n\geq1}(-1)^n\,\frac{2z}{z^2 - n^2} .
\]
Verificação de sinais: $\operatorname{Res}\bigl(\pi/\sin(\pi z),
n\bigr) = \pi/(\pi\cos\pi n) = (-1)^n$
(\cref{met:b3:residues:compute}), e $(-1)^n\frac{2z}{z^2 -
n^2} = \frac{(-1)^n}{z-n} + \frac{(-1)^n}{z+n}$ carrega
exatamente esse \omterm{def:b3:residues:singularities}{resíduo} em $\pm n$.

\textbf{21.} Para $\abs z \leq r < 1$, expandindo cada termo como
na questão 19 ($(-1)^n\frac{2z}{z^2-n^2} =
(-1)^{n-1}\frac{2z}{n^2}\sum_k\frac{z^{2k}}{n^{2k}}$) e
aplicando Fubini:
\[
\frac{\pi}{\sin(\pi z)} = \frac1z +
2\sum_{k\geq0}\eta(2k+2)\,z^{2k+1},
\qquad
\eta(s) = \sum_{n\geq1}\frac{(-1)^{n-1}}{n^s} .
\]
Lado de Taylor: $\sin(\pi z) = \pi z(1 - \frac{(\pi z)^2}6 +
O(z^4))$ dá $\frac\pi{\sin\pi z} = \frac1z +
\frac{\pi^2}6z + O(z^3)$: $2\eta(2) = \frac{\pi^2}6$, de modo que
$\eta(2) = \frac{\pi^2}{12}$. Coerência: $\eta(2) =
\zeta(2) - 2\sum_n(2n)^{-2} = (1 - 2^{1-2})\zeta(2) =
\frac{\zeta(2)}2 = \frac{\pi^2}{12}$.

\textbf{22.} Em $z = \frac12$, a questão 20 dá
\[
\pi = 2 + \sum_{n\geq1}(-1)^n\frac{1}{\frac14 - n^2}
= 2 + 4\sum_{n\geq1}\frac{(-1)^{n-1}}{(2n-1)(2n+1)} .
\]
Com $\frac1{(2n-1)(2n+1)} = \frac12\bigl(\frac1{2n-1} -
\frac1{2n+1}\bigr)$, e sendo ambas as séries alternadas convergentes
(critério de Leibniz), a soma se reparte como $\frac12\bigl[L -
(1 - L)\bigr] = L - \frac12$, em que $L = 1 - \frac13 +
\frac15 - \cdots$ e $\frac13 - \frac15 + \frac17 - \cdots
= 1 - L$. Logo $\pi = 2 + 4(L - \frac12) = 4L$:
\[
\frac\pi4 = 1 - \frac13 + \frac15 - \frac17 + \cdots
\]
A moral: cada núcleo é \omterm{def:b3:residues:singularities}{meromorfo} com polos numa
progressão aritmética e \omterm{def:b3:residues:singularities}{resíduos} prescritos. A
cotangente põe \omterm{def:b3:residues:singularities}{resíduo} $1$ em todo inteiro e soma
$f(n)$; sua derivada eleva os polos ao quadrado e precifica
$\lambda(2)$; a tangente move os polos para $\frac12 + \Z$
e precifica os denominadores ímpares; $\pi/\sin$ mantém os
polos inteiros, mas alterna os \omterm{def:b3:residues:singularities}{resíduos} $(-1)^n$, daí
as séries alternadas. Os núcleos de primeira ordem são todos funções
ímpares: emparelhar $n$ com $-n$ duplica os coeficientes de potência par
e aniquila os ímpares, de modo que $\zeta(2k)$
jorra mecanicamente, ao passo que $\zeta(3)$ nunca aparece. A
cegueira é de paridade, não falta de técnica.

\textbf{23.} Com $w = 2\iu\pi z$:
\[
\iu\pi z + \frac{2\iu\pi z}{\eu^{2\iu\pi z} - 1}
= \iu\pi z\,\frac{\eu^{2\iu\pi z} + 1}{\eu^{2\iu\pi z} - 1}
= \iu\pi z\,\frac{\eu^{\iu\pi z} +
\eu^{-\iu\pi z}}{\eu^{\iu\pi z} - \eu^{-\iu\pi z}}
= \pi z\,\frac{\cos\pi z}{\sin\pi z} = \pi z\cot(\pi z) .
\]
Logo, usando a função geradora em $w = 2\iu\pi z$ (e
$B_1 = -\frac12$ cancelando o termo $\iu\pi z$, anulando-se os $B$ ímpares
além disso):
\[
\pi z\cot(\pi z) = \sum_{k\geq0}\frac{B_{2k}}{(2k)!}
(2\iu\pi z)^{2k}
= 1 + \sum_{k\geq1}(-1)^k\frac{(2\pi)^{2k}B_{2k}}{(2k)!}
z^{2k} .
\]
Por outro lado, a expansão em frações parciais (Parte IV)
dá $\pi z\cot(\pi z) = 1 - 2\sum_{k\geq1}\zeta(2k)z^{2k}$
(expanda cada $\frac{2z^2}{z^2 - n^2} =
-2\sum_k\frac{z^{2k}}{n^{2k}}$ e some em $n$, justificando a convergência
normal a troca para $\abs z < 1$).
Comparando coeficientes: $-2\zeta(2k) =
(-1)^k\frac{(2\pi)^{2k}B_{2k}}{(2k)!}$, a fórmula enunciada.

\textbf{24.} $k = 1$: $\frac{(2\pi)^2}{2\cdot2}\cdot\frac16
= \frac{\pi^2}6$. $k = 2$: $-\frac{(2\pi)^4}{2\cdot24}
\cdot\bigl(-\frac1{30}\bigr) = \frac{16\pi^4}{48\cdot30} =
\frac{\pi^4}{90}$. $k = 3$: $\frac{(2\pi)^6}{2\cdot720}
\cdot\frac1{42} = \frac{64\pi^6}{1440\cdot42} =
\frac{\pi^6}{945}$.

\textbf{25.} Escreva $C(z) = \pi z\cot(\pi z) = 1 -
2\sum_{k\geq1}\zeta(2k)z^{2k}$ (questão 23). Derivação
direta de $C = \pi z\cot(\pi z)$, usando $(\cot u)'
= -1 - \cot^2u$:
\[
zC'(z) = \pi z\cot(\pi z) - \pi^2z^2\bigl(1 +
\cot^2(\pi z)\bigr) = C - \pi^2z^2 - C^2 .
\]
Agora expanda ambos os membros em potências de $z^2$. Membro esquerdo:
$\sum_k(-4k)\,\zeta(2k)\,z^{2k}$. Membro direito: elevando a série ao
quadrado,
\[
C - C^2 = 2\sum_{k\geq1}\zeta(2k)z^{2k}
- 4\sum_{k\geq2}\Bigl(\sum_{j=1}^{k-1}\zeta(2j)
\zeta(2k-2j)\Bigr)z^{2k},
\]
e o termo $-\pi^2z^2 = -6\zeta(2)z^2$ apenas ajusta $k =
1$. Comparando os coeficientes de $z^{2k}$ para $k \geq 2$:
\[
-4k\,\zeta(2k) = 2\,\zeta(2k) -
4\sum_{j=1}^{k-1}\zeta(2j)\,\zeta(2k-2j),
\]
isto é, $\bigl(k + \frac12\bigr)\zeta(2k) =
\sum_{j=1}^{k-1}\zeta(2j)\zeta(2k-2j)$. (Em $k = 1$ a
identidade se lê $-4\zeta(2) = 2\zeta(2) - 6\zeta(2)$: uma
verificação de coerência, não informação nova.) Aplicações: $k =
2$: $\frac52\zeta(4) = \zeta(2)^2 = \frac{\pi^4}{36}$, de modo que
$\zeta(4) = \frac{\pi^4}{90}$; $k = 3$: $\frac72\zeta(6) =
2\zeta(2)\zeta(4) = \frac{\pi^6}{270}$, de modo que $\zeta(6) =
\frac{2}{7}\cdot\frac{\pi^6}{270} = \frac{\pi^6}{945}$. Uma
semente transcendente ($\zeta(2) = \frac{\pi^2}6$), e a álgebra pura
gera todo valor par da zeta.
\end{solution}
